% ============================================================================
% Fichier  : partie2/P02_C06_gouvernance.tex
% Titre    : Gouvernance de la Sécurité
% Auteur   : MINKA MI NGUIDJOÏ Thierry Emmanuel
% Version  : 1.0 -- novembre 2026
% ============================================================================
% Base : brouillon « Livre II » (gouvernance, politiques, rôles) et « Livre
% VIII » (sécurité administrative). Sources : ISO/IEC 27014:2020 ; ISO/IEC
% 27001:2022 (clause 5) ; ISO/IEC 27004:2016 ; loi n° 2010/012, art. 26 (1),
% texte officiel.
% ============================================================================

\chapter{Gouvernance de la Sécurité}
\label{chap:gouvernance}
\chaptermark{Gouvernance}

\epigraph{%
    « La sécurité se décide en haut,\\
    se met en œuvre partout, et se vérifie toujours. »%
}{Principe du cours}

\niveauChapitre{Fondamental}{Chapitres~\ref{chap:proprietes}
et~\ref{chap:menaces}.}

\begin{boxobjectifs}
\begin{itemize}
    \item Distinguer gouvernance et management de la sécurité.
    \item Répartir les rôles et les responsabilités, et les formaliser dans
          une matrice.
    \item Construire une politique de sécurité et la documentation qui
          en découle.
    \item Définir des indicateurs de sécurité utiles à la direction.
    \item Intégrer la sécurité aux projets et aux contrats.
\end{itemize}
\end{boxobjectifs}

\bigskip

% ============================================================================
\section*{Question directrice : qui était responsable de la sécurité de LiZbiZ~?}
% ============================================================================

Personne. M.~KOUNOU gérait l'informatique~; SAHEL INFOGÉRANCE exploitait
deux centres~; la direction pensait que l'antivirus et les pare-feu
« faisaient la sécurité ». Aucune politique, aucun responsable, aucun
indicateur. L'Opération GAROUA n'est pas seulement une faille technique~:
c'est une absence de gouvernance.

\medskip
\textbf{Plan du chapitre.}
\begin{enumerate}
    \item Gouverner et manager la sécurité (section~\ref{sec:c06-gouv}).
    \item Les rôles et les responsabilités (section~\ref{sec:c06-roles}).
    \item La politique de sécurité (section~\ref{sec:c06-pssi}).
    \item Les indicateurs (section~\ref{sec:c06-indicateurs}).
    \item La sécurité dans les projets et les contrats (section~\ref{sec:c06-projets}).
    \item Une gouvernance pour LiZbiZ (section~\ref{sec:c06-lizbiz}).
\end{enumerate}

\separateur

% ============================================================================
\section{Gouverner et Manager la Sécurité}
\label{sec:c06-gouv}
% ============================================================================

\begin{boxdef}
\termecle{Gouvernance de la sécurité de l'information}~: système par
lequel la direction oriente et contrôle les activités de sécurité~: elle
fixe les objectifs, alloue les ressources, accepte les risques et
vérifie les résultats. Le \termecle{management de la sécurité} met en
œuvre ces orientations au quotidien~\cite{iso27014_2020}.
\end{boxdef}

L'ISO/IEC~27014 décrit cinq processus de gouvernance~: \emph{évaluer}
(la situation et les risques), \emph{orienter} (fixer la politique et les
priorités), \emph{surveiller} (suivre les résultats), \emph{communiquer}
(aux parties intéressées) et \emph{assurer} (faire vérifier de façon
indépendante, par un audit).

\begin{boxjuris}
\textbf{\loicam{n°~2010/012}, article~26, alinéa~1} (texte officiel)~:
\textit{« Les exploitants des systèmes d'information prennent toutes les
mesures techniques et administratives afin de garantir la sécurité des
services offerts. A cet effet, ils se dotent de systèmes normalisés leur
permettant d'identifier, d'évaluer, de traiter et de gérer continûment les
risques liés à la sécurité des systèmes d'information dans le cadre des
services offerts directement ou indirectement. »}

\medskip
\textbf{Analyse.} La loi exige un dispositif \emph{continu} et
\emph{normalisé} de gestion des risques~: c'est une obligation de
gouvernance, pas seulement de moyens techniques. Un exploitant qui n'a ni
responsable, ni appréciation des risques, ni suivi, comme LiZbiZ avant
mars, ne peut pas démontrer qu'il la respecte.
\end{boxjuris}

% ============================================================================
\section{Les Rôles et les Responsabilités}
\label{sec:c06-roles}
% ============================================================================

\begin{tableaucours}{L{3.4cm}L{8.8cm}}{%
    \tableauHeader{Acteur} &
    \tableauHeader{Responsabilités}
}
Direction générale & Approuve la politique, fixe l'appétence au risque,
accepte les risques résiduels, alloue le budget \\
Comité de sécurité & Réunit direction, métiers, informatique et RSSI~;
arbitre les priorités, suit le plan \\
RSSI & Propose la politique, apprécie les risques, pilote le plan,
conseille, contrôle, rend compte \\
Propriétaires d'actifs & Responsables métier d'une information ou d'un
processus~; fixent ses besoins de sécurité et valident les droits \\
Direction informatique & Met en œuvre les mesures techniques~;
exploite \\
Utilisateurs & Respectent les règles~; signalent les incidents \\
Prestataires & Appliquent les exigences du contrat~; rendent compte \\
Audit & Vérifie, de façon indépendante, que le dispositif fonctionne \\
\end{tableaucours}
\vspace{-0.8em}
\begin{center}{\small\itshape Tableau~6.1. Les acteurs de la sécurité}\end{center}

\subsection{Le rattachement du RSSI}

Un RSSI rattaché à la direction informatique contrôle son propre
supérieur~: il a du mal à signaler qu'un correctif n'est pas installé. Un
rattachement à la direction générale, ou à une fonction de gestion des
risques, lui donne l'indépendance nécessaire. C'est un choix de
gouvernance, que la direction motive par écrit.

\subsection{La matrice RACI}

\begin{tableaucours}{L{4.2cm}cccc}{%
    \tableauHeader{Activité} &
    \tableauHeader{DG} &
    \tableauHeader{RSSI} &
    \tableauHeader{DSI} &
    \tableauHeader{Propriétaire}
}
Approuver la politique & A & R & C & C \\
Apprécier les risques & I & A & C & R \\
Installer les correctifs & I & C & A/R & I \\
Valider les droits d'accès & I & C & R & A \\
Accepter un risque résiduel & A & R & C & C \\
Gérer un incident majeur & A & R & R & C \\
\end{tableaucours}
\vspace{-0.8em}
\begin{center}{\small\itshape Tableau~6.2. Matrice RACI (R réalise, A
approuve et rend compte, C consulté, I informé)}\end{center}

% ============================================================================
\section{La Politique de Sécurité}
\label{sec:c06-pssi}
% ============================================================================

\begin{boxdef}
\termecle{Politique de sécurité des systèmes d'information} (PSSI)~:
document approuvé par la direction qui énonce les objectifs de sécurité,
les principes et les responsabilités, et que déclinent des politiques
thématiques et des procédures. L'ISO/IEC~27001 en fait une exigence de la
direction (clause~5.2)~\cite{iso27001_2022}.
\end{boxdef}

\begin{tableaucours}{L{3.6cm}L{5.0cm}L{3.6cm}}{%
    \tableauHeader{Niveau} &
    \tableauHeader{Contenu} &
    \tableauHeader{Approuvé par}
}
Politique générale & Objectifs, principes, rôles, engagement & Direction
générale \\
Politiques thématiques & Accès, mots de passe, sauvegardes, télétravail,
prestataires, développement & Comité de sécurité \\
Standards & Règles techniques précises (configurations, longueurs de clés) &
RSSI \\
Procédures & Pas à pas d'une opération (créer un compte, restaurer) &
Responsable du service \\
Guides & Recommandations non obligatoires & RSSI \\
\end{tableaucours}
\vspace{-0.8em}
\begin{center}{\small\itshape Tableau~6.3. La pyramide documentaire}\end{center}

\begin{boxalert}
\textbf{Violation typique.} Copier la politique d'une autre organisation,
la faire signer, et la ranger.

\smallskip
\textbf{Règle~:} une politique part des risques de l'organisation
(Chapitre~\ref{chap:risques}), est connue de ceux qui l'appliquent, et
est revue au moins chaque année et après chaque incident majeur. Une
politique que personne n'applique est une preuve contre l'organisation.
\end{boxalert}

\begin{boxprocedure}
\textbf{Faire vivre la politique}
\begin{enumerate}
    \item \textbf{Rédiger} à partir de l'appréciation des risques.
    \item \textbf{Faire approuver} par la direction, avec une date de revue.
    \item \textbf{Diffuser} à chaque personne concernée, et le tracer.
    \item \textbf{Former} aux règles qui les concernent.
    \item \textbf{Contrôler} l'application, par des indicateurs et des
          audits.
    \item \textbf{Réviser} après chaque incident et chaque changement
          majeur.
\end{enumerate}
\end{boxprocedure}

% ============================================================================
\section{Les Indicateurs}
\label{sec:c06-indicateurs}
% ============================================================================

L'ISO/IEC~27004 donne une méthode pour mesurer l'efficacité du dispositif
de sécurité~\cite{iso27004_2016}. Un bon indicateur a une définition, une
formule, une source, une fréquence, un responsable et une cible.

\begin{tableaucours}{L{4.2cm}L{4.4cm}L{3.6cm}}{%
    \tableauHeader{Indicateur} &
    \tableauHeader{Formule} &
    \tableauHeader{Cible (exemple)}
}
Correctifs critiques à jour & Serveurs à jour / serveurs, à 30~jours &
100~\% \\
Comptes distants avec double authentification & Comptes protégés /
comptes distants & 100~\% \\
Restaurations testées & Tests réussis / tests prévus, par trimestre &
100~\% \\
Délai de détection & Médiane entre début d'un incident et sa détection &
Moins de 24~h \\
Personnes formées & Formées dans l'année / effectif & Plus de 90~\% \\
\end{tableaucours}
\vspace{-0.8em}
\begin{center}{\small\itshape Tableau~6.4. Indicateurs de sécurité pour une
direction}\end{center}

\begin{boxinfo}
\textbf{Mesurer aussi ce que l'on ignore.} Un indicateur dont la source
n'existe pas (aucun journal, aucun inventaire) ne vaut pas zéro~: il vaut
« inconnu ». Le Chapitre~\ref{chap:mesurer} montre comment le dire
(\NEK{}) plutôt que de le taire.
\end{boxinfo}

% ============================================================================
\section{La Sécurité dans les Projets et les Contrats}
\label{sec:c06-projets}
% ============================================================================

La sécurité coûte moins cher quand elle est prévue. Chaque projet
informatique inclut une analyse de risques, des exigences de sécurité,
des tests avant mise en production et une validation par le RSSI. Chaque
contrat avec un prestataire qui touche au système d'information contient
des clauses de sécurité.

\begin{tableaucours}{L{4.0cm}L{8.2cm}}{%
    \tableauHeader{Clause} &
    \tableauHeader{Contenu}
}
Exigences de sécurité & Double authentification, correctifs, comptes
nominatifs \\
Journaux & Conservation et remise des journaux d'accès sur demande \\
Notification des incidents & Délai (par exemple 24~h) et contenu \\
Droit d'audit & Possibilité d'auditer ou de faire auditer le prestataire \\
Sous-traitance & Autorisation préalable de tout sous-traitant \\
Réversibilité & Restitution des données et des configurations en fin de
contrat \\
\end{tableaucours}
\vspace{-0.8em}
\begin{center}{\small\itshape Tableau~6.5. Clauses de sécurité d'un contrat
de prestation}\end{center}

% ============================================================================
\section{Une Gouvernance pour LiZbiZ}
\label{sec:c06-lizbiz}
% ============================================================================

\begin{boxscenario}{Opération GAROUA. Les décisions du conseil d'avril}
\textbf{Décisions.} Le conseil d'administration de LiZbiZ décide~:
(1)~de recruter un RSSI rattaché au directeur général~; (2)~de créer un
comité de sécurité trimestriel~; (3)~d'adopter une politique générale et
cinq politiques thématiques avant septembre~; (4)~de renégocier le contrat
de SAHEL INFOGÉRANCE avec les clauses du Tableau~6.5~; (5)~de présenter
chaque trimestre cinq indicateurs au conseil.

\medskip
\textbf{Analyse} (confiance moyenne). Ces décisions répondent aux causes
d'organisation de l'attaque. Leur efficacité ne se jugera qu'aux
indicateurs, dans six mois~: une décision n'est pas un résultat.
\end{boxscenario}

\begin{boxmetier}{RSSI}
\textbf{Parler à la direction.} Le RSSI présente des risques en langage
métier (jours sans facturation, données exposées, sanctions), propose
des options chiffrées et laisse la direction décider. Il trace par écrit
chaque risque accepté.
\end{boxmetier}

\begin{boxmetier}{Auditeur de sécurité}
\textbf{Vérifier la gouvernance.} L'auditeur demande la politique signée
et sa date de revue, la nomination du RSSI, les comptes rendus du comité,
les indicateurs et les risques acceptés. Une politique sans preuve de
diffusion ni de contrôle se cote comme absente de fait.
\end{boxmetier}

\separateur

\begin{boxresume}
\textbf{Ce qu'il faut retenir de ce chapitre~:}
\begin{itemize}
    \item \textbf{La gouvernance} évalue, oriente, surveille, communique
          et fait assurer~; le management met en œuvre.
    \item \textbf{La loi n°~2010/012} (art.~26) impose une gestion continue
          et normalisée des risques de sécurité.
    \item \textbf{Les rôles} se formalisent (matrice RACI)~; le RSSI a
          besoin d'indépendance.
    \item \textbf{La politique} s'appuie sur les risques, se diffuse, se
          contrôle et se révise~; elle se décline en politiques thématiques,
          standards et procédures.
    \item \textbf{Les indicateurs} ont une formule, une source, une cible~;
          une source absente donne « inconnu », pas zéro.
    \item \textbf{Les contrats} portent des clauses de sécurité.
\end{itemize}
\end{boxresume}

\begin{boxinfo}
\textbf{Points de contrôle de l'élève.} Je sais~: citer les cinq
processus de gouvernance de l'ISO/IEC~27014~; remplir une matrice RACI~;
décrire la pyramide documentaire~; rédiger un indicateur complet~; lister
six clauses de sécurité d'un contrat.
\end{boxinfo}

\bigskip

\begin{boxexo}[title={\ding{45}\quad Exercice~6.1. Gouverner ou manager~? \niveaubase}]
Classez~: approuver la politique~; installer un pare-feu~; accepter un
risque résiduel~; former les utilisateurs~; fixer le budget de sécurité~;
revoir les droits d'accès.
\end{boxexo}

\begin{boxexo}[title={\ding{45}\quad Exercice~6.2. Le texte de l'article 26 \niveaubase}]
Relevez, dans l'article~26, les quatre verbes qui décrivent la gestion des
risques, et associez à chacun une activité concrète.
\end{boxexo}

\begin{boxexo}[title={\ding{45}\quad Exercice~6.3. RACI \niveaumoyen}]
Complétez la matrice RACI pour~: créer un compte~; supprimer un compte
d'un salarié parti~; restaurer une sauvegarde~; choisir un prestataire.
\end{boxexo}

\begin{boxexo}[title={\ding{45}\quad Exercice~6.4. Rattachement \niveaumoyen}]
La BANQUE SANAGA hésite à rattacher son RSSI au DSI ou au directeur des
risques. Donnez deux arguments pour chaque option et votre recommandation.
\end{boxexo}

\begin{boxexo}[title={\ding{45}\quad Exercice~6.5. Pyramide \niveaubase}]
Classez dans la pyramide documentaire~: « les mots de passe comptent au
moins 14~caractères »~; « LiZbiZ protège les données de ses clients »~;
« pour restaurer une base, suivre les étapes 1 à 9 »~; « politique de
télétravail ».
\end{boxexo}

\begin{boxexo}[title={\ding{45}\quad Exercice~6.6. Un indicateur \niveaumoyen}]
Rédigez la fiche complète de l'indicateur « délai d'installation des
correctifs critiques »~: définition, formule, source, fréquence,
responsable, cible.
\end{boxexo}

\begin{boxexo}[title={\ding{45}\quad Exercice~6.7. Le contrat \niveaumoyen}]
Pour chaque clause du Tableau~6.5, dites quel maillon de l'Opération
GAROUA elle aurait touché.
\end{boxexo}

\begin{boxexo}[title={\ding{45}\quad Exercice~6.8. La politique copiée \niveaumoyen}]
Un auditeur trouve chez l'ANPP une politique de 60~pages, copiée d'une
banque européenne, signée en 2019 et jamais revue. Rédigez la
constatation et la recommandation.
\end{boxexo}

\begin{boxexo}[title={\ding{45}\quad Exercice~6.9. Le tableau de bord \niveauexpert}]
Proposez au conseil de LiZbiZ un tableau de bord trimestriel de six
indicateurs, et dites pour chacun ce que LiZbiZ pourrait mesurer
aujourd'hui, et ce qui serait « inconnu ».
\end{boxexo}

\begin{boxexo}[title={\ding{45}\quad Exercice~6.10. La politique générale \niveauexpert}]
Rédigez la politique générale de sécurité de LiZbiZ en une page~: objet,
champ, objectifs, principes, rôles, révision.
\end{boxexo}

\bigskip

\begin{boxinfo}
\textbf{Pour aller plus loin.}
\begin{itemize}
    \item ISO/IEC~27014, gouvernance de la sécurité de
          l'information~\cite{iso27014_2020}.
    \item ISO/IEC~27004, mesure de la sécurité~\cite{iso27004_2016}.
\end{itemize}
\end{boxinfo}

\begin{boxinfo}
\textbf{Transition.} La gouvernance décide sur la base des risques. Le
Chapitre~\ref{chap:risques} montre comment les apprécier et les traiter.
\end{boxinfo}
